\chapter{Conclusion}
\label{ch:essay-conclusion}

The four theorems examined in this essay establish that representation, coherence, repair, and reconstruction are not four independent subjects requiring separate theoretical foundations but four aspects of a single admissibility relation examined at different points in its operation. Representation Invariance addresses the question of when physically distinct substrates count as semantically equivalent, answered by checking whether the admissibility structures they induce are related by an appropriate isomorphism. Admissibility Preservation addresses the question of what makes a trajectory coherent, answered by verifying that the reachable set remains nontrivial at each step rather than degenerating to emptiness or a predetermined singleton. Monotonic Relaxation addresses the question of when inadmissibility decreases, answered conditionally for trajectories following repair dynamics that increase reachable-set volume. Non-Identifiability addresses the question of when reconstruction becomes impossible in principle, answered by showing that observation maps factoring through reachable-set structure are generically non-injective whenever history space exceeds what that structure encodes.

What unifies these four results is that each derives from the same underlying apparatus: the admissibility relation $x \adm{R} y$ determining which continuations remain admissible from any given state, and the reachable set $\reach{x}$ collecting those admissible continuations into a structure that can be examined, compared, and projected. The formal development in \emph{Admissible Reconstruction} establishes that no additional primitive vocabulary is required beyond this minimal starting point. Representation equivalence is admissibility-structure isomorphism. Coherence is reachable-set nondegeneracy. Repair is reachable-set volume increase. Non-identifiability is equivalence-class structure under projection. The economy is not coincidental but reflects the fact that a single relation, when examined from different perspectives, generates all four phenomena as consequences.

The cross-modal investigation developed here applies this unified apparatus to the projection problem: determining what structural relations survive when latent intentional states are realized through different physical channels. Speech, gesture, music, and audio differ in which aspects of structure each modality preserves and which it necessarily discards, but the formal question remains the same across modalities. Does the projection preserve admissibility structure completely, as required by Theorem~1.1 for representational equivalence, or does it introduce distortions that break the isomorphism? Does the projected trajectory maintain nontrivial reachable sets at each step, satisfying the coherence condition of Theorem~2.1, or does it degenerate? If the system undergoes repair, does the trajectory follow the volume-increasing dynamics required by Theorem~3.1 for monotonic relaxation? And when an observer attempts to reconstruct latent structure from observable projection, does the observation map produce collisions as predicted by Theorem~4.2, and if so, are those collisions correctable or structural as diagnosed by Corollary~4.1?

The deepest result concerns the limits of reconstruction. Non-identifiability establishes that when distinct histories produce the same observable outcomes, the failure to distinguish them is not a weakness of the observer's inferential methods but a property of the observation channel itself. Two histories that share the same reachable-set structure have become, from the perspective of any observation map preserving only continuation structure, identical in the observable projection. The information required to distinguish them has not merely become difficult to access; it has ceased to exist in the projection. This is what it means for observables to be lossy: not that some details are omitted while others remain, as though careful observation could in principle recover the missing information, but that entire dimensions of variation in the history space collapse to equivalence classes in the observable space, and no refinement of observational technique operating within that space can recover the collapsed dimensions.

The diagnostic corollary provides the criterion for recognizing when such collapse has occurred. Given two candidate histories producing similar observables, check whether their reachable sets coincide. If they differ, the reconstruction uncertainty is correctable, and additional evidence could in principle resolve it. If they are equal, the uncertainty is structural, and no observation map factoring through reachable-set structure can eliminate it. This transforms reconstruction failure from an undifferentiated acknowledgment that evidence may be insufficient into a structured diagnosis with decidable answers: the failure is either correctable or fundamental, and the distinction can be determined by examining the reachable-set structure directly.

The framework developed here therefore addresses the projection problem not by asking whether reconstruction succeeds or fails in particular cases but by establishing what structural properties determine success and failure. Representations are equivalent when admissibility-isomorphisms exist. Trajectories are coherent when reachable sets remain nontrivial. Repair succeeds when reachable-set volume increases. Reconstruction is impossible when distinct histories share reachable-set structure. These are not four separate claims requiring independent justification but four views of one relation, and recognizing them as such clarifies both what the admissibility apparatus establishes and what it leaves for empirical investigation to determine.
